% ID: SEL-D0012-BD-E08
% ESTADO: cotejada. Segunda pasada visual completa y borrador compilado revisados.
% Fuente: Selectividad/matematicas_aplicadas/EXTRA_RESERVA_B_Modelo-6-Enunciados (Mat-Ap-2021).pdf
% SHA-256: 1752ec2f46b4ef2cc5d1c3346d69b79c2cc8a7aa12f6a0cd8f0d994b5bcda157
% Página: 2; bloque D; ejercicio original 8.

La estatura de las mujeres de una población sigue una ley Normal de media
desconocida y desviación típica 7 cm.
\begin{enumerate}[label=\alph*)]
  \item \textbf{(1.5 puntos)} Se toma una muestra aleatoria de 300 mujeres
  de esta población, que da una estatura media de 168 cm. Construya un
  intervalo de confianza al 97\% para estimar la estatura media de las
  mujeres de esta población.
  \item \textbf{(1 punto)} Calcule el tamaño mínimo que debe tener una muestra
  de esta población para que, con un nivel de confianza del 94\%, el error
  máximo cometido al estimar la estatura media de las mujeres de esa población
  sea inferior a 1.2 cm.
\end{enumerate}
